\section{Competence: is it actually good at this, or only good on paper?}
\label{sec:competence}

The first question a clinician asks of any tool is whether it is competent at the
specific task, on patients like the ones in this clinic. Competence is not a single
benchmark number; it is a stack of evidence built in order, from bench accuracy to
prospective validation on the site's own population. The convergence of human
judgment and machine pattern recognition only earns its place at the bedside when
the machine has been shown to perform, not in the abstract, but on the question and
the population in front of the clinician \cite{topol2019high}.

This is where the mechanism of the framework begins to matter. Every action the
system contemplates is produced under verification before generation: Claude Code
proposes a candidate step, Codex reviews that step independently, and a ten-gate
VVUQ examination (verification that the action was built correctly, validation that
it is the right action for this patient, and uncertainty quantification of how
confident the system is) resolves the proposal to one of three outcomes: ACCEPT
and act under audit, ESCALATE to a qualified human, or BLOCK before the patient.
Competence, in this framing, is not a marketing claim but the documented record
that this examination has been passed on evidence the clinician can inspect.

The author's assurance lineage shows what that record looks like in practice. A
surgical-humanoid platform built for a mobile pancreatic-cancer unit was subjected
to a priority VVUQ examination before any action was permitted
\cite{Kawchak2026MobilePancreaticCancerUnitree}, and the broader verification-
automation work demonstrates how the same gate suite is applied to an oncology
clinical trial rather than to a single procedure
\cite{Kawchak2026VVUQOncologyClinicalTrial}. The headline number from the assurance
run is concrete and checkable: the system passed 172 of 172 automated tests before
it was cleared to act. A clinician does not have to take competence on faith; the
clinician can ask to see which tests were run and how each resolved, which is the
discipline the illustrative H. R. 9510 would write into the trial record
\cite{Kawchak2026HR9510Bill}.

The decisive adoption action here is local validation. A model that performs well
on a vendor's bench, or even on a published retrospective cohort, has not yet been
shown to perform on this site's patients, with this site's imaging protocols and
case mix. Competence is therefore earned on a ladder, and each rung answers a
different question while leaving others open. \autoref{tab:competence} sets out the
ladder so the clinician can see precisely what a given level of evidence
establishes and, just as important, what it cannot. The competence evidence stack
in Figure 3 traces the same ascent, ending only when the assurance run closes the
gap between paper performance and the population in the room.

\begin{cfwfig}
\node[cfone] (a) at (0,0) {Bench}; \node[cftwo] (b) at (3.0,0) {Retrospective};
\node[cftwo] (c) at (6.4,0) {Prospective}; \node[cfact] (d) at (9.8,0) {Assurance 172/172};
\node[cfhope] (e) at (6.4,-2.0) {Competence established};
\draw[cfedge] (a)--(b); \draw[cfedge] (b)--(c); \draw[cfedge] (c)--(d);
\draw[cfedge] (d) to[out=-90,in=0] (e.east);
\end{cfwfig}
\vspace{-0.2cm}
\cfwcaption{Figure 3. The competence evidence stack (Mermaid 06 reproduced as
colored TikZ): bench accuracy, retrospective cohort, prospective on-site
validation, and the assurance run advance in order until competence is established
on the local population.}

\needspace{9\baselineskip}\begin{table}[H]
\footnotesize
\begin{tabularx}{\textwidth}{@{}L{3cm} Y L{3.4cm}@{}}
\toprule
\textbf{Evidence level} & \textbf{What it establishes} & \textbf{What it cannot show}\\
\midrule
Bench accuracy & The model meets a numeric performance threshold on a curated reference set & That it holds on real patients or this site's protocols\\
Retrospective cohort & Performance on historical records spanning a real population & That it acts safely in real time, before an outcome is known\\
Prospective on-site validation & Performance on this clinic's own patients under the actual workflow & That rare or catastrophic cases are handled, which require the gates\\
Assurance run (172 of 172 gate tests) & Every verification, validation, and uncertainty gate passed before action was permitted & A substitute for continued monitoring once the system is in routine use\\
\bottomrule
\end{tabularx}
\caption{The competence evidence ladder: what each level of evidence establishes
and what it leaves unanswered, culminating in local validation and the assurance
run.}
\label{tab:competence}
\end{table}

Competence is answered not by a claim but by evidence the clinician can inspect,
and by the discipline of validating on the local population before clinical use.
The honest test is the family test: would I let this system near my own family
member on the strength of a bench number alone, or only after I had seen it
validated on patients like the one in front of me and watched it pass every gate?
